% ============================================================
% Section 6 — Analysis & Ablations
% ============================================================
\section{Analysis}
\label{sec:analysis}

\subsection{Ablation Study}
\label{sec:analysis:ablation}

To understand the contribution of each component, we ablate key modules of HAT using the Phi-3-3.8B Cortex on the NQ benchmark.
Results are summarised in Table~\ref{tab:ablation}.

\begin{table}[t]
  \centering
  \caption{Ablation study on the NQ benchmark (Phi-3-3.8B).  Each row removes or replaces one component.  $\Delta$ indicates the accuracy change relative to the full HAT system.}
  \label{tab:ablation}
  \small
  \begin{tabular}{@{}lcc@{}}
    \toprule
    \textbf{Variant} & \textbf{Accuracy} & $\bm{\Delta}$ \\
    \midrule
    \textbf{HAT (full)}                      & -- & ---    \\
    \midrule
    \quad w/o Abstraction                     & -- & --     \\
    \quad w/o Selection (random write)        & -- & --     \\
    \quad w/o Novelty signal ($\gamma = 0$)   & -- & --     \\
    \quad w/o Uncertainty signal ($\alpha = 0$)& -- & --    \\
    \quad w/o Feedback signal ($\beta = 0$)   & -- & --     \\
    \quad w/o Oracle consultation             & -- & --     \\
    \quad w/o EWC regularisation              & -- & --     \\
    \quad w/o Replay Builder (raw traces)     & -- & --     \\
    \bottomrule
  \end{tabular}
\end{table}

\paragraph{Abstraction is critical.}
% TODO: Discuss impact of removing the abstraction stage.

\paragraph{Selection outperforms random curation.}
% TODO: Discuss random selection vs. write-policy-based selection.

\paragraph{All three selection signals contribute.}
% TODO: Discuss individual contribution of uncertainty, feedback, novelty.

\paragraph{EWC prevents forgetting.}
% TODO: Discuss forgetting rate with and without EWC.

\subsection{Effect of SWS Cycle Frequency}
\label{sec:analysis:frequency}

We vary the number of interactions between SWS cycles ($N \in \{100, 250, 500, 1000, 2000\}$) and measure final accuracy and computational cost.

\begin{figure}[t]
  \centering
  % TODO: Replace with actual plot
  \fbox{\parbox{0.46\textwidth}{\centering\vspace{4em}\textbf{[SWS Frequency Plot Placeholder]}\vspace{4em}}}
  \hfill
  \fbox{\parbox{0.46\textwidth}{\centering\vspace{4em}\textbf{[Forgetting Rate Plot Placeholder]}\vspace{4em}}}
  \caption{\textbf{Left:} Final accuracy as a function of SWS cycle frequency.  \textbf{Right:} Forgetting rate across SWS cycles for HAT vs.\ baselines.}
  \label{fig:sws_frequency}
\end{figure}

% TODO: Discuss optimal cycle frequency and trade-offs.

\subsection{Qualitative Analysis}
\label{sec:analysis:qualitative}

\paragraph{What does the Hippocampus select?}
We analyse the distribution of memory traces written to the Neocortex across interaction types.
% TODO: Add breakdown figure or table.
% Expected: user corrections and high-uncertainty responses dominate, while routine interactions are filtered out.

\paragraph{Memory evolution over time.}
We visualise the Neocortex content across SWS cycles using t-SNE embeddings.
% TODO: Add t-SNE figure showing how the memory distribution shifts.
% Expected: early cycles accumulate broad knowledge; later cycles focus on edge cases and user-specific preferences.

\paragraph{Case study: error correction propagation.}
We trace a specific user correction through the HAT pipeline—from raw interaction to abstracted trace, selection, replay construction, and its effect on model behaviour after SWS training.
% TODO: Add illustrative example.

\subsection{Scaling Behaviour}
\label{sec:analysis:scaling}

We examine how HAT's gains vary with Cortex model size (1.5B, 3.8B, 8B).
% TODO: Discuss whether smaller models benefit more from HAT (expected: yes, since they have more room for improvement and stronger models have less uncertainty).
